\documentclass[11pt,a4paper]{article}
\usepackage[T1]{fontenc}
\usepackage{amsmath,amssymb}
\usepackage{booktabs}
\usepackage{hyperref}
\pdfstringdefDisableCommands{%
  \def\pi{pi}%
  \def\zeta{zeta}%
  \def\mathrm#1{#1}%
  \def\textbf#1{#1}%
}
\usepackage{xcolor}
\usepackage{tcolorbox}
\usepackage[margin=22mm]{geometry}
\usepackage{xeCJK}
\setCJKmainfont{Hiragino Mincho ProN}
\setCJKsansfont{Hiragino Sans}
\tcbuselibrary{breakable}

% Custom boxes
\newtcolorbox{tipbox}[1][]{
  colback=yellow!5, colframe=orange!60!black,
  fonttitle=\bfseries, title=コツ, breakable, #1}
\newtcolorbox{warnbox}[1][]{
  colback=red!5, colframe=red!60!black,
  fonttitle=\bfseries, title=警告, breakable, #1}
\newtcolorbox{candidatebox}[1][]{
  colback=blue!3, colframe=blue!50!black,
  fonttitle=\bfseries, title={#1}, breakable}

\title{Setup B' 実装ガイドブック\\[6pt]
\large 室温 MEMS + DBR による Boyer 符号反転実験の完全マニュアル\\[4pt]
\normalsize 独立研究者・アマチュア向け、共同研究先アプローチ戦略付き}
\author{Wright Brothers Research Program}
\date{2026年4月11日}

\begin{document}
\maketitle

\begin{abstract}
本ガイドブックは、Stage 1（Boyer 符号反転）の Setup B'（室温 MEMS + DBR）を
\textbf{独立研究者・アマチュア}が実際に走らせるための完全マニュアルである。

実験の本質は「標準カシミール力測定装置の片方の電極を、Akkermans 公式で
逆設計した DBR 多層膜に置き換え、力の符号と大きさを比較する」こと。
予算はシナリオ A（共同研究前提）で約 ¥1M、期間は 6 ヶ月。

本書の中核は \textbf{候補共同研究先の詳細プロフィール}と
\textbf{独立研究者からの具体的アプローチ戦略}である。各候補者を
「アマチュアへのオープンさ」「アプローチのしやすさ」の評価軸で
1--5 のスコアでランク付けし、それぞれに最適化された連絡方法・
提案書テンプレート・初回メール例を提供する。

成功への鍵は「反重力」を前面に出さず、\textbf{「Boyer (1974) の 50 年未検証の
予言を、Akkermans 公式によるオイラー積的解釈で初検証する」}という
古典的に正しく、知的に興味深い切り口で提案することである。
\end{abstract}

\tableofcontents
\newpage

%======================================================================
\section{はじめに：このガイドの読者と目的}
%======================================================================

\subsection{想定読者}

\begin{itemize}
\item 独立研究者・アマチュア物理愛好家で、自前の実験設備を持たない人
\item 大学・研究所の物理学者と \textbf{共同研究} を立ち上げたい人
\item 既存の Casimir 力測定装置を活用して \textbf{新しい実験} を提案したい人
\item 「アカデミアの常識」を理解せずにいきなりアプローチして失敗したくない人
\end{itemize}

\subsection{このガイドが扱わないこと}

\begin{itemize}
\item Casimir 力の基礎物理（既存の教科書を参照、例：Milonni \emph{The Quantum Vacuum}）
\item 装置の詳細製作技術（共同研究先の専門領域）
\item WB 反重力理論の詳細（\texttt{topological\_antigravity.tex} を参照）
\end{itemize}

\subsection{重要な前提}

\textbf{Setup B' は反重力を直接検証する実験ではない}。
本実験は古典電磁気学の枠内で完結する Boyer (1974) の 50 年未検証予言を
DBR 多層膜で初実験的に確かめるものである。

\begin{tipbox}
共同研究先にアプローチするとき「反重力」「ワープドライブ」といった用語を
\textbf{絶対に最初に出してはいけない}。Boyer 効果の初検証としてピッチすれば、
古典物理として完全に正当な提案となり、受け入れられる確率が桁違いに上がる。
反重力との関係は、共同研究が軌道に乗ってから（あるいは論文投稿後に）
\emph{副次的な}話題として触れる程度に留める。
\end{tipbox}

\newpage

%======================================================================
\section{実験の概要：30 秒で}
%======================================================================

\subsection{提案する実験}

標準の MEMS / AFM 系カシミール力測定装置を用い、片方の電極を：

\begin{itemize}
\item \textbf{制御群}：Au コート Si プレート（標準的な金属電極）
\item \textbf{処理群}：SiO$_2$/Ta$_2$O$_5$ 多層 DBR
  （Akkermans 公式で逆設計、DN 等価境界条件を実装）
\end{itemize}

の 2 種類で交換し、カシミール力 $F(L)$ を $L = 100$ nm$\sim$1 µm の範囲で
測定して比較する。

\subsection{予測}

理論予測（Boyer 1974, Kenneth--Klich 2002 で厳密化）：
\[
  \frac{F_{\rm DBR}}{F_{\rm Au}} = -\frac{7}{8} = -0.875
\]
すなわち：
\begin{itemize}
\item 符号反転（引力 $\to$ 斥力）
\item 大きさは Au-Au の $7/8$ 倍
\end{itemize}

\subsection{50 年未検証の理由}

Boyer 1974 が予言した PEC--PMC（完全電気導体--完全磁気導体）境界条件での
カシミール斥力は、PMC が天然材料として存在しないために誰も実現したことがない。
本実験は \textbf{DBR 多層膜が PMC 等価境界条件を実装する}という新しい
洞察（Akkermans 公式 + 数論的逆設計）を使うことで、この 50 年来の問題を
解く。

\subsection{独立した科学的価値}

たとえ反重力理論が間違っていても、本実験単体で：

\begin{itemize}
\item Boyer 50 年未検証予言の初検証
\item Akkermans 公式の物理応用としての初実装
\item PMC 等価境界の新規工学的実現
\end{itemize}

これらは \textbf{Nature / PRL クラスの論文}になりうる単独価値を持つ。

\newpage

%======================================================================
\section{理論的根拠（要点のみ）}
%======================================================================

\subsection{Boyer ratio の導出}

理想 PEC 板間（DD 境界）：
\[
  E_{\rm DD}/A = -\frac{\pi^2 \hbar c}{720 L^3}\;,\qquad F_{\rm DD}/A = -\frac{\pi^2 \hbar c}{240 L^4}
\]

PEC--PMC 板間（DN 境界）：モードが整数 $n$ から半整数 $(n-1/2)$ に変わる。
これにより：
\[
  \frac{F_{\rm DN}}{F_{\rm DD}} = \frac{-\zeta(-3,1/2)}{-\zeta(-3)}
                                = \frac{+7/960}{-1/120} = -\frac{7}{8}
\]

ここで $\zeta(-3) = 1/120$ は Riemann ゼータ、$\zeta(-3, 1/2) = -7/960$ は
Hurwitz ゼータ（Bernoulli 多項式 $B_4(1/2)$ から計算）。

\subsection{DBR が PMC 等価になる理由：Akkermans 公式}

Akkermans--Dunne--Levy (2013) の式：
\[
  \zeta_H(s) - \zeta_{H_0}(s) = \frac{1}{\pi}\int_0^\infty dE\,E^{-s}\,\frac{d\delta(E)}{dE}
\]

これは「1D 多層膜の S 行列位相シフト $\delta(E)$ から、その構造のスペクトル
ゼータが完全に決まる」という関係。逆に解けば、目標スペクトル（DN 等価）を
実装するための位相シフトが計算でき、それを再現する多層膜層厚を inverse design
できる。

\subsection{設計仕様（参考値）}

\begin{itemize}
\item 材料：SiO$_2$ ($n_L = 1.45$) / Ta$_2$O$_5$ ($n_H = 2.10$)
\item 層数：15--30 ペア
\item Bragg 中心波長 $\lambda_0$：実験距離 $L$ に対して $\lambda_0 \approx 2L \times \sqrt{\varepsilon_{\rm eff}}$
  程度（$L = 100$ nm なら $\lambda_0 \sim 280$ nm 近辺、UV-A 帯）
\item 設計目標：Bragg 中心で $|r|^2 \approx 1$、$\arg(r) \approx 0$（PMC 等価）
\item 詳細は \texttt{notes\_2026\_04\_10\_phase2\_audit\_ja.pdf} §7 を参照
\end{itemize}

\newpage

%======================================================================
\section{必要な装置・材料リスト}
%======================================================================

\subsection{A: カシミール力測定装置（最大項目）}

\begin{center}
\begin{tabular}{p{0.20\textwidth}p{0.45\textwidth}p{0.25\textwidth}}
\toprule
オプション & 内容 & コスト \\
\midrule
A1 既存 AFM 流用 & 共用施設の Bruker / JPK / Asylum AFM、カスタム球プローブ装着 & ¥0--¥200k（プローブのみ） \\
A2 MEMS トーション秤 & Decca-IUPUI 系、MEMS ファブで製作 & ¥500k--¥2M \\
A3 商用 AFM 新規購入 & Bruker Dimension Icon 等 & ¥10--15M \\
\bottomrule
\end{tabular}
\end{center}

\textbf{独立研究者向け推奨：A1}。共同研究先の AFM を借用する。

\subsection{B: DBR 多層膜（実験の核心）}

\begin{itemize}
\item \textbf{材料}：SiO$_2$/Ta$_2$O$_5$ または TiO$_2$/SiO$_2$
\item \textbf{表面粗さ}：$< 1$ nm RMS（ion-beam sputtering 必須）
\item \textbf{基板}：fused silica または Si、$\lambda/10$ 平坦度
\item \textbf{調達先}：
  \begin{itemize}
  \item \textbf{Layertec GmbH}（独）：custom DBR、¥300--500k、納期 6--8 週
  \item \textbf{Newport / MKS}：standard supermirror、¥200--400k
  \item \textbf{Advanced Thin Films}（米）：IBS、¥500k--1M
  \item \textbf{LIGO/KAGRA レガシー}：重力波検出器の余剰品、共同研究で ¥0
  \end{itemize}
\end{itemize}

\subsection{C: 参照プレート}

Au コート Si or fused silica、表面粗さ $< 1$ nm RMS、Au 厚 200 nm。
コスト：¥50--100k。

\subsection{D--G: 周辺装置}

\begin{center}
\begin{tabular}{lll}
\toprule
項目 & 仕様 & 中古コスト \\
\midrule
D 真空チャンバー    & $10^{-6}$ Torr、$\sim 30$ cm 径、turbo + scroll & ¥200--500k \\
E 振動絶縁          & 空圧除振台 (Newport / Thorlabs)              & ¥500k or 借用 \\
F1 ピエゾステージ    & Z 軸 nm 精度（PI / Aerotech）                 & ¥500k--1M \\
F2 Kelvin 補正回路   & DC バイアス源 + 制御                          & ¥150k \\
G ロックインアンプ   & SR830 / Zurich MFLI                           & ¥300k--1M \\
\bottomrule
\end{tabular}
\end{center}

シナリオ A では D--G は \textbf{全て共同研究先の既存設備を借用}する。

\newpage

%======================================================================
\section{予算 3 シナリオ}
%======================================================================

\subsection{シナリオ A：最安構成（既存ラボとの共同研究、推奨）}

\begin{center}
\begin{tabular}{lr}
\toprule
項目 & 金額 \\
\midrule
DBR 多層膜（Layertec custom） & ¥400k \\
Au 球プローブ + 参照 Au プレート & ¥150k \\
Kelvin 補正回路（自作 or 既製） & ¥100k \\
消耗品（ガス、配線、洗浄液） & ¥100k \\
旅費・滞在費（3 月、海外想定） & ¥300k \\
\midrule
\textbf{合計} & \textbf{¥1.05M} \\
\bottomrule
\end{tabular}
\end{center}

\textbf{前提}：AFM、真空チャンバー、ピエゾ、ロックイン、振動絶縁を全て共同研究先で借用。

\subsection{シナリオ B：自前構築（中規模、参考）}

\begin{center}
\begin{tabular}{lr}
\toprule
項目 & 金額 \\
\midrule
DBR 多層膜 & ¥500k \\
カスタム MEMS or 高性能 AFM プローブ & ¥1M \\
中古真空チャンバー + ポンプ & ¥500k \\
ピエゾステージ & ¥800k \\
ロックインアンプ（中古） & ¥300k \\
振動絶縁テーブル & ¥500k \\
Kelvin 補正系 & ¥150k \\
消耗品 & ¥200k \\
人件費（6 月） & ¥1M \\
\midrule
\textbf{合計} & \textbf{¥4.95M} \\
\bottomrule
\end{tabular}
\end{center}

\subsection{シナリオ C：完全新規（参考、非推奨）}

シナリオ B の装置を全て新品で揃えた場合：\textbf{合計約 ¥24M}。
独立研究者には現実的でない。

\subsection{推奨選択}

\textbf{シナリオ A 一択}。本実験の本質的価値は装置ではなく DBR 設計と
力測定の再解釈にあるので、装置投資は無駄。共同研究先の選定がすべて。

\newpage

%======================================================================
\section{候補共同研究先の詳細プロフィール（本ガイドの中核）}
%======================================================================

\subsection{評価軸について}

各候補者を以下 5 軸で 1--5 スコアで評価する。

\begin{center}
\begin{tabular}{p{0.30\textwidth}p{0.62\textwidth}}
\toprule
評価軸 & 意味 \\
\midrule
\textbf{科学的適合性} & その人の研究内容が Setup B' に直接関係するか \\
\textbf{装置適合性} & 必要な実験設備（AFM/MEMS, DBR 評価系）を持っているか \\
\textbf{オープンさ} & アマチュア・非伝統的研究への態度 \\
\textbf{アプローチ容易性} & 連絡しやすさ、返信率、ゲートキーピングの少なさ \\
\textbf{成功時の影響力} & その人と組んだ場合の論文の影響力・可視性 \\
\bottomrule
\end{tabular}
\end{center}

\textbf{重要な注意}：以下の評価は \emph{公開情報}（発表論文、所属機関の広報、
公開講演、起業活動等）から推測したものである。各個人の実際の性向・スタンスは
直接接触して確認するしかない。\textbf{これは絶対視すべき格付けではなく、
最初のアプローチ順序を決めるための参考に過ぎない}。

\vspace{1em}

\subsection{総合ランキング（独立研究者向け）}

\begin{center}
\renewcommand{\arraystretch}{1.4}
\begin{tabular}{cllc}
\toprule
順位 & 候補者 & 所属 & 総合 \\
\midrule
\textbf{1} & Davide Iannuzzi      & Vrije Univ.\ Amsterdam / Optics11 (蘭) & ★★★★★ \\
2 & Alexander O.\ Sushkov & Boston Univ. (米)         & ★★★★$\frac{1}{2}$ \\
3 & Jeremy N.\ Munday    & UC Davis (米)              & ★★★★ \\
4 & Steve K.\ Lamoreaux  & Yale Univ. (米)            & ★★★$\frac{1}{2}$ \\
5 & Ho Bun Chan          & HKUST (香港)               & ★★★ \\
6 & Umar Mohideen        & UC Riverside (米)          & ★★★ \\
7 & Ricardo S.\ Decca    & IUPUI (米)                 & ★★$\frac{1}{2}$ \\
\midrule
8 & 国内候補：理研、NTT 物性研、産総研 & --- & ★★ (個別調査要) \\
\bottomrule
\end{tabular}
\end{center}

\textbf{独立研究者には上位 3 名を強く推奨}。以下、各候補者の詳細プロフィール。

\newpage

%======================================================================
\section{候補 1：Davide Iannuzzi（最有力）}
%======================================================================

\begin{candidatebox}[Davide Iannuzzi]
\textbf{所属}：Vrije Universiteit Amsterdam（VU Amsterdam）、
LaserLaB 物理学部教授。同時に \textbf{Optics11} の共同創業者・CEO。

\textbf{Casimir 関連論文}：
\begin{itemize}
\item \textbf{Iannuzzi, Lisanti, Capasso (2004) PNAS 101, 4019}：
  ``Effect of hydrogen-switchable mirrors on the Casimir force.''
  \textbf{鏡の状態変化でカシミール力を変調する初実験}。これがまさに本実験の
  系譜。
\item Iannuzzi 系の MEMS Casimir 力測定論文多数
\end{itemize}

\textbf{現在の研究}：ファイバー光学センサー、医療応用、Casimir は副業的継続
\end{candidatebox}

\subsection{評価}

\begin{center}
\begin{tabular}{lc}
\toprule
評価軸 & スコア \\
\midrule
科学的適合性    & ★★★★★ \\
装置適合性      & ★★★★★ \\
オープンさ      & ★★★★★ \\
アプローチ容易性 & ★★★★ \\
成功時の影響力  & ★★★★ \\
\midrule
\textbf{総合}   & \textbf{★★★★★} \\
\bottomrule
\end{tabular}
\end{center}

\subsection{なぜ最有力か}

\paragraph{1. 直接の系譜.}
Iannuzzi 2004 (PNAS) は \emph{まさに本実験の前駆}：「鏡の状態を変えてカシミール
力を変調する」という発想を実装した最初の論文。Setup B' は彼の 2004 論文の
\emph{拡張} として完全に位置付けられる。提案を受けたとき、彼は 30 秒で
本質を理解する。

\paragraph{2. 起業家マインド.}
Optics11 の共同創業者・CEO。ファイバー光センサー製品を世界中の研究機関に
販売している。\textbf{学術ゲートキーピングが薄く、外部からの提案を機会とみる
傾向}が強い。アマチュア起業家の発想にも開かれている可能性が高い。

\paragraph{3. 産業化志向.}
``Boyer 50 年予言の初検証'' という具体的な成果に直結する提案は、
産業化志向の研究者にとって魅力的（PRL/Nature 級の単独成果）。

\paragraph{4. 装置完備.}
VU Amsterdam の Casimir 系装置は世界トップクラスで、AFM、MEMS、真空系、
振動絶縁を全て持っている。

\subsection{アプローチ戦略：3 ステップ}

\paragraph{Step 1：論文を準備する.}
arXiv にプレプリントを 1 本投稿する。タイトル例：

\begin{quote}
\itshape
``Akkermans inverse design of artificial magnetic conductors for the
experimental verification of Boyer's 50-year-old prediction of repulsive
Casimir forces''
\end{quote}

内容は本ガイドの §3 + §7（Akkermans 公式 + DBR 設計）を 6--8 ページにまとめる。
Wright Brothers 名義で出す。これは \textbf{calling card} になる。

\paragraph{Step 2：個人メールを送る.}
Iannuzzi の VU Amsterdam の公開メールアドレスに、簡潔な個人メールを送る。
テンプレートは §15 にある。重要なポイント：

\begin{itemize}
\item \textbf{冒頭で 2004 PNAS 論文を引用}し、それの自然な拡張だと位置付ける
\item arXiv 番号を入れる
\item ``反重力'' は一切書かない
\item \textbf{DBR は自分が手配する}（彼にコストを負わせない）
\item ``あなたの研究室を 2 週間訪問させてもらえないか''と聞く
\end{itemize}

\paragraph{Step 3：返信があったら 30 分のオンライン会議.}
オンライン会議で：
\begin{itemize}
\item Akkermans 解析の数値結果（PDF 1 枚）を画面共有
\item DBR の Layertec 見積もりを提示
\item 「測定は完全にあなたの装置と手順で。私は DBR と労力を提供する」
\item 共著論文の \textbf{第一著者を Iannuzzi にする}と明言
\end{itemize}

\subsection{成功確率の自己評価}

公開情報からの主観的見積もり：\textbf{30--50\%}。
これは独立研究者からの提案としては \emph{極めて高い}。

理由：(1) 2004 論文の直接拡張、(2) Optics11 経由の起業家マインド、
(3) 提案者がコストを負担、(4) 結果が PRL/Nature 級。

\textbf{失敗するとしたら}：本人が現在ファイバーセンサー事業で多忙すぎて
アカデミック実験への時間が取れない場合。その場合は彼の研究室の博士課程
学生・ポスドクに紹介してもらう道がある。

\newpage

%======================================================================
\section{候補 2：Alexander O.\ Sushkov}
%======================================================================

\begin{candidatebox}[Alexander Sushkov]
\textbf{所属}：Boston University 物理学部 Associate Professor。

\textbf{Casimir 関連論文}：
\begin{itemize}
\item \textbf{Sushkov, Kim, Dalvit, Lamoreaux (2011) Nature Phys.\ 7, 230}：
  ``Observation of the thermal Casimir force.''
  熱カシミール効果の初観測。
\item 各種カシミール精密測定論文
\end{itemize}

\textbf{他の研究}：
\begin{itemize}
\item \textbf{CASPEr（Cosmic Axion Spin Precession Experiment）}：
  超軽量 axion ダークマター探索。NMR + 精密磁場測定。
\item 原子物理、精密測定、基礎物理一般
\end{itemize}
\end{candidatebox}

\subsection{評価}

\begin{center}
\begin{tabular}{lc}
\toprule
評価軸 & スコア \\
\midrule
科学的適合性    & ★★★★ \\
装置適合性      & ★★★★ \\
オープンさ      & ★★★★★ \\
アプローチ容易性 & ★★★★ \\
成功時の影響力  & ★★★★$\frac{1}{2}$ \\
\midrule
\textbf{総合}   & \textbf{★★★★$\frac{1}{2}$} \\
\bottomrule
\end{tabular}
\end{center}

\subsection{なぜ第 2 候補か}

\paragraph{1. 「fringe topics を真剣にやる」で有名.}
Sushkov は \textbf{未検出の axion ダークマター} という超ハイリスク・
超ハイリターン研究に長年コミットしている。これは \emph{ほとんどの主流
物理学者がやらない種類の研究}で、彼が新規・非伝統的提案に対して
\textbf{構造的にオープン}であることを示している。

\paragraph{2. Lamoreaux の弟子.}
Sushkov 2011 Nature Phys.\ の共著者は Lamoreaux（候補 4）。Lamoreaux 系列の
Casimir 文化を継承しつつ、より広い基礎物理に展開している。Boyer 50 年予言
という \textbf{古典的に正当だが未検証}の課題は、彼の研究スタイルに完全に合致。

\paragraph{3. 若い世代.}
Decca や Mohideen より若く、新しい実験提案に時間を割く余裕がある。

\subsection{アプローチ戦略}

\paragraph{Step 1：CASPEr との接続を作る.}
最初のメールで「Boyer 効果の精密測定は、CASPEr で培ったあなたの基礎物理測定の
専門知識に直結する」と書く。彼の主要研究との相補性を強調する。

\paragraph{Step 2：論文 + Akkermans 計算結果.}
arXiv 論文を添えるのは Iannuzzi の場合と同じ。加えて、Akkermans 公式の
数値解の Python コードを GitHub で公開して、リンクを提供する。
Sushkov は \textbf{コードを書く実験家}なので、再現可能な計算は強い説得材料。

\paragraph{Step 3：``リスク共有''の提案.}
「私は DBR と Layertec 見積もりを持参します。実験が失敗しても、Akkermans
公式の物理応用としての論文化は可能です（JPhys A など）。リスクを最小化
するための提案です」と明示する。

\subsection{成功確率}

\textbf{20--35\%}。Iannuzzi より低いのは、Sushkov の主戦場が axion で、
Casimir は副業的だから。しかし「提案を理解する能力」と「アマチュアへの
寛容さ」は最高クラス。

\newpage

%======================================================================
\section{候補 3：Jeremy N.\ Munday}
%======================================================================

\begin{candidatebox}[Jeremy Munday]
\textbf{所属}：UC Davis 電気・コンピュータ工学部教授。

\textbf{Casimir 関連論文}：
\begin{itemize}
\item \textbf{Munday, Capasso, Parsegian (2009) Nature 457, 170}：
  ``Measured long-range repulsive Casimir-Lifshitz forces.''
  \textbf{斥力カシミール力の初測定}（液体中の誘電率コントラストによる）。
\item Casimir 力の精密測定多数（Capasso 系の系譜）
\end{itemize}

\textbf{現在の研究}：太陽電池、ナノフォトニクス、放射冷却。Casimir は継続的副業。
\end{candidatebox}

\subsection{評価}

\begin{center}
\begin{tabular}{lc}
\toprule
評価軸 & スコア \\
\midrule
科学的適合性    & ★★★★$\frac{1}{2}$ \\
装置適合性      & ★★★★ \\
オープンさ      & ★★★$\frac{1}{2}$ \\
アプローチ容易性 & ★★★$\frac{1}{2}$ \\
成功時の影響力  & ★★★★ \\
\midrule
\textbf{総合}   & \textbf{★★★★} \\
\bottomrule
\end{tabular}
\end{center}

\subsection{なぜ第 3 候補か}

\paragraph{1. 唯一の「斥力カシミール初測定」の実施者.}
Munday 2009 Nature は \emph{現在のところ唯一の斥力カシミールの実験報告}で、
Setup B' の提案を最も深く理解できる人物。\textbf{ただし機構は誘電率コントラストで、
DN 境界ではない}ので、Setup B' は彼の 2009 論文の \emph{別経路の補完}として
位置付けられる。

\paragraph{2. Capasso 系の出身.}
Federico Capasso（Harvard）系列。Capasso 自身よりは若く、アクセスしやすい。

\paragraph{3. 工学部所属.}
工学部の人は応用志向が強く、「具体的に何ができるか」を重視する傾向。
実装可能な提案に乗りやすい。

\subsection{アプローチ戦略}

\paragraph{Step 1：``Munday 2009 を超えて''というフレーミング.}
\begin{quote}
\itshape
``Your 2009 Nature paper showed repulsive Casimir via dielectric contrast,
but the original Boyer prediction (PEC--PMC) has never been verified.
We propose to verify Boyer's specific mechanism using Akkermans-designed DBR.''
\end{quote}
彼自身の論文を出発点として尊重し、その「未踏領域」に進む提案だと位置付ける。

\paragraph{Step 2：応用面を強調.}
工学部所属なので、純粋な基礎物理よりも \textbf{応用可能性} を強調する：
``DN 等価境界の制御は将来的にナノフォトニクス・カシミール工学に応用可能''。

\paragraph{Step 3：UC Davis の学生にコンタクト経路を作る.}
Munday 本人が忙しいなら、彼の学生・ポスドクに先にアプローチする。
研究室サイトで現メンバーを確認し、似た研究をしている人にメール。

\subsection{成功確率}

\textbf{15--30\%}。Iannuzzi と Sushkov より低いのは、(1) 工学部志向が強く
基礎物理 50 年予言の検証への熱意が読みにくい、(2) 現在の主戦場（太陽電池等）
からカシミールへの再注目が必要だから。

\newpage

%======================================================================
\section{候補 4：Steve K.\ Lamoreaux}
%======================================================================

\begin{candidatebox}[Steve Lamoreaux]
\textbf{所属}：Yale University 物理学部教授。

\textbf{Casimir 関連論文}：
\begin{itemize}
\item \textbf{Lamoreaux (1997) PRL 78, 5}：
  カシミール力の現代的初測定。歴史的論文。
\end{itemize}

\textbf{他の研究}：axion ダークマター探索（ADMX 系）、hidden photon、
中性子の電気双極子モーメント探索など、\textbf{基礎物理探索の重鎮}。
\end{candidatebox}

\subsection{評価}

\begin{center}
\begin{tabular}{lc}
\toprule
評価軸 & スコア \\
\midrule
科学的適合性    & ★★★★ \\
装置適合性      & ★★★ \\
オープンさ      & ★★★★ \\
アプローチ容易性 & ★★ \\
成功時の影響力  & ★★★★★ \\
\midrule
\textbf{総合}   & \textbf{★★★$\frac{1}{2}$} \\
\bottomrule
\end{tabular}
\end{center}

\subsection{なぜ第 4 候補か}

\paragraph{1. 歴史的名声.}
Lamoreaux の名前が論文に入れば PRL クラスは確実、Nature も視野。
\emph{現代カシミール物理の創始者} の一人としての権威。

\paragraph{2. 基礎物理のフロンティア常駐.}
axion・hidden photon・EDM 等、\textbf{未確認・非主流の基礎物理を真剣にやる}
ことで知られる。Boyer 50 年予言の検証は彼の興味の中心軸。

\paragraph{3. ただし忙しい・高名.}
これが弱点。連絡しても返信が来ない可能性が高い。学生・ポスドクのフィルタを
通る必要がある。

\subsection{アプローチ戦略}

\paragraph{Step 1：Sushkov 経由で間接接触.}
Lamoreaux への直接アプローチより、\textbf{Sushkov（候補 2）経由でのリファラル}
を狙う。Sushkov は Lamoreaux の弟子で、Sushkov に提案が響けば
Lamoreaux に紹介してくれる可能性がある。

\paragraph{Step 2：会議で接触.}
APS March Meeting や DAMOP などの会議で、彼の発表後の Q\&A セッションで
1 つ良い質問をする。これがフックになる。

\paragraph{Step 3：直接メール（最後の手段）.}
メール本文は短く、目的を 3 文以内に圧縮する：
\begin{quote}
\itshape
``Boyer 1974's PEC--PMC repulsive Casimir prediction has never been
experimentally verified. We have an Akkermans-based DBR design that
implements the equivalent boundary condition. Would you have 30 minutes
for a Zoom call to discuss feasibility?''
\end{quote}
arXiv 論文へのリンクのみ添付。長いメールは捨てられる。

\subsection{成功確率}

\textbf{5--15\%}。Iannuzzi の半分以下。ただし \emph{成功した場合の影響力} は最大。
独立研究者の最初の選択肢としては推奨しない（他の候補と並行して試す）。

\newpage

%======================================================================
\section{候補 5：Ho Bun Chan}
%======================================================================

\begin{candidatebox}[Ho Bun Chan]
\textbf{所属}：Hong Kong University of Science and Technology (HKUST)
物理学部教授。

\textbf{Casimir 関連論文}：
\begin{itemize}
\item \textbf{Tang, Wang, Rodriguez, Reid, Johnson, Chan (2017) Nat.\ Photonics
  11, 97}：``Measurement of non-monotonic Casimir forces between silicon
  nanostructures.'' \textbf{幾何効果による斥力 Casimir の測定}。
\item 各種 MEMS Casimir 測定論文
\end{itemize}
\end{candidatebox}

\subsection{評価}

\begin{center}
\begin{tabular}{lc}
\toprule
評価軸 & スコア \\
\midrule
科学的適合性    & ★★★$\frac{1}{2}$ \\
装置適合性      & ★★★★ \\
オープンさ      & ★★★ \\
アプローチ容易性 & ★★★ \\
成功時の影響力  & ★★★ \\
\midrule
\textbf{総合}   & \textbf{★★★} \\
\bottomrule
\end{tabular}
\end{center}

\subsection{なぜ候補か}

幾何構造による斥力 Casimir を測定した \textbf{Munday 2009 と並ぶ貴重な実験家}。
また、HKUST はアジアの研究機関なので、欧米の Casimir コミュニティとは
やや独立した文化を持ち、新規提案に対する反応が読みにくい（オープンの可能性も
ガチガチの可能性もある）。

\subsection{アプローチ戦略}

英語メールで Iannuzzi と同じテンプレートを使う。違いは：
\begin{itemize}
\item Tang 2017 を冒頭で引用
\item ``幾何効果と DN 境界効果は補完的''と位置付ける
\item HKUST 訪問の費用負担を明確に提示（自己負担）
\end{itemize}

\subsection{成功確率}

\textbf{10--20\%}。

\newpage

%======================================================================
\section{候補 6：Umar Mohideen}
%======================================================================

\begin{candidatebox}[Umar Mohideen]
\textbf{所属}：UC Riverside 物理・天文学部教授。

\textbf{Casimir 関連論文}：
\begin{itemize}
\item \textbf{Mohideen, Roy (1998) PRL 81, 4549}：
  AFM ベースの精密 Casimir 測定。Lamoreaux 1997 と並ぶ初期の重要論文。
\item Klimchitskaya, Mostepanenko らとの精密測定論文多数
\end{itemize}
\end{candidatebox}

\subsection{評価}

\begin{center}
\begin{tabular}{lc}
\toprule
評価軸 & スコア \\
\midrule
科学的適合性    & ★★★★ \\
装置適合性      & ★★★★★ \\
オープンさ      & ★★$\frac{1}{2}$ \\
アプローチ容易性 & ★★ \\
成功時の影響力  & ★★★★ \\
\midrule
\textbf{総合}   & \textbf{★★★} \\
\bottomrule
\end{tabular}
\end{center}

\subsection{プロファイル}

カシミール測定の最高峰の一人。\textbf{装置と精度では世界トップ}。しかし
Klimchitskaya--Mostepanenko 系の理論的な細部にこだわる文化に属し、
\textbf{新規・非伝統的な提案には保守的}な可能性がある。

\textbf{推奨度}：他の候補が全て失敗した場合の保険。最初に attack しない。

\subsection{アプローチ戦略}

\paragraph{Step 1：彼の最近の論文を精読する.}
Mohideen と Klimchitskaya--Mostepanenko の最近 5 年分の論文を読み、
彼らが何を「正しい」「間違っている」と主張しているか把握する。彼らの理論的
立場と矛盾しない形で提案を構成する。

\paragraph{Step 2：理論的に厳密な提案書.}
Iannuzzi よりも \textbf{遥かに厳密で長い}提案書を準備する。
Akkermans 公式の導出、誤差バジェット、系統誤差解析を全て付ける。
``アマチュア''風の簡潔な提案では捨てられる可能性が高い。

\paragraph{Step 3：Klimchitskaya 経由を狙う.}
Mohideen 本人ではなく、共同研究者の Galina Klimchitskaya（理論家、ロシア）に
先にコンタクトする方法もある。彼女が興味を示せば Mohideen にもつながる。

\subsection{成功確率}

\textbf{5--15\%}。

\newpage

%======================================================================
\section{候補 7：Ricardo S.\ Decca}
%======================================================================

\begin{candidatebox}[Ricardo Decca]
\textbf{所属}：Indiana University--Purdue University Indianapolis (IUPUI)
物理学部教授。

\textbf{Casimir 関連}：
\begin{itemize}
\item MEMS torsion balance を用いた sub-pN 級の精密 Casimir 測定
\item Klimchitskaya--Mostepanenko 系の主要実験家
\item Drude vs plasma model 論争の主要プレイヤー
\end{itemize}
\end{candidatebox}

\subsection{評価}

\begin{center}
\begin{tabular}{lc}
\toprule
評価軸 & スコア \\
\midrule
科学的適合性    & ★★★★ \\
装置適合性      & ★★★★★ \\
オープンさ      & ★$\frac{1}{2}$ \\
アプローチ容易性 & ★$\frac{1}{2}$ \\
成功時の影響力  & ★★★$\frac{1}{2}$ \\
\midrule
\textbf{総合}   & \textbf{★★$\frac{1}{2}$} \\
\bottomrule
\end{tabular}
\end{center}

\subsection{プロファイル}

\textbf{装置と測定技術では世界最高峰}。MEMS torsion balance の決定版。
ただし Mohideen 系列と同じく Klimchitskaya--Mostepanenko 文化に属し、
\textbf{独立研究者からの提案に対する開放度は低い可能性が高い}。

\textbf{推奨度}：独立研究者の最初の選択肢としては不適。
ただし、Iannuzzi/Sushkov 経由で彼に紹介されれば話は別。

\subsection{アプローチ戦略}

\textbf{直接アプローチは推奨しない}。他の候補が成功して論文が出た後、
論文の追試・拡張の文脈で接触するのが現実的。

\newpage

%======================================================================
\section{国内候補（理研、NTT、産総研）}
%======================================================================

\subsection{総評}

日本国内の Casimir 力測定を直接の専門とするグループは少ない。
ただし \textbf{隣接領域}（MEMS、超伝導 cQED、ナノフォトニクス）には
強いグループが存在する。

\subsection{候補機関}

\paragraph{理研 CEMS（創発物性科学研究センター）.}
\begin{itemize}
\item 量子物理・物性理論の総合拠点
\item 中村泰信（量子情報）系列で SQUID 系の経路（Setup B''）には強い
\item Setup B' の AFM カシミール測定には直接の専門家は少ないかもしれない
\item アプローチ：Web 公開の連絡先からメール。日本語可。
\end{itemize}

\paragraph{NTT 物性科学基礎研究所（厚木）.}
\begin{itemize}
\item MEMS、超伝導回路、ナノファブの強力なインフラ
\item 産業研究所なので外部との共同研究には独自のプロセスがある
\item アプローチ：技術提案として正式に申し込む
\end{itemize}

\paragraph{産業技術総合研究所（産総研）.}
\begin{itemize}
\item 計測標準研究部門が精密測定の専門家を抱える
\item Casimir 測定の前例があるかは要調査
\item アプローチ：産総研の「技術相談」窓口経由が比較的開かれている
\end{itemize}

\subsection{独立研究者にとっての国内候補のメリット・デメリット}

\paragraph{メリット.}
\begin{itemize}
\item 言語障壁ゼロ
\item 訪問・打ち合わせコストが低い
\item 同じタイムゾーン
\end{itemize}

\paragraph{デメリット.}
\begin{itemize}
\item 国内の Casimir 専門家層が薄い（欧米と比較）
\item 産業研究所はアマチュア提案に冷淡な傾向
\item 大学は所属確認・信用情報を求められる場合が多い
\end{itemize}

\textbf{推奨}：海外候補（特に Iannuzzi）と並行してアプローチする「保険」として
扱う。

\newpage

%======================================================================
\section{アマチュア $\to$ アカデミアへのアプローチ：一般的なコツ}
%======================================================================

\subsection{大原則}

\begin{tipbox}
\textbf{アカデミア研究者は ``proof of seriousness'' を求めている}。
あなたが時間を投資して理解した証拠と、あなた自身がコストを負担する意思を
示せば、内容が良ければ多くの研究者は反応する。
\end{tipbox}

\subsection{やってはいけない 7 つの NG}

\begin{enumerate}
\item \textbf{``反重力''や ``ワープドライブ'' を冒頭に書く}
  $\to$ 即座に pseudoscience フォルダ行き
\item \textbf{``私の理論を検証してください''}
  $\to$ 押しつけと受け取られる。``共同で'' を強調する
\item \textbf{長すぎるメール（500 単語超）}
  $\to$ 読まずに削除される
\item \textbf{メールに添付ファイル多数}
  $\to$ ウイルスと混同される。リンクで提供
\item \textbf{所属詐称}
  $\to$ 即発覚、信用ゼロ
\item \textbf{``Nobel 級''などの誇張表現}
  $\to$ 即座に怪しい人扱い
\item \textbf{返信を急かす}
  $\to$ 返信率が下がる。1 ヶ月待つ
\end{enumerate}

\subsection{逆に効果的な 7 つの行動}

\begin{enumerate}
\item \textbf{arXiv プレプリントを先に出す}
  $\to$ 物理学者は arXiv 番号で「真面目さ」を即座に評価する
\item \textbf{相手の論文を 1 つ正確に引用}
  $\to$ 相手の研究を読んでいる証拠
\item \textbf{具体的な提案（装置、予算、期間、役割分担）}
  $\to$ 「話を聞くだけ」と違い真剣味が伝わる
\item \textbf{自己負担の明示（DBR は私が手配します）}
  $\to$ 相手のリスクを下げる
\item \textbf{論文の第一著者を相手にする}と最初に明言
  $\to$ アカデミア通貨である「論文業績」を相手に渡す
\item \textbf{オンライン会議の打診（30 分）}
  $\to$ 対面より圧倒的に承諾率が高い
\item \textbf{失敗時の代替論文化経路を用意}
  $\to$ ``失敗しても何か論文になる''は研究者を安心させる
\end{enumerate}

\subsection{独立研究者特有の壁と回避策}

\paragraph{壁 1：所属がない.}
回避：``Independent Researcher'' と書く。pseudonymity（Wright Brothers）は
GitHub プロフィールと arXiv プレプリントが揃っていれば許容される。
``Bourbaki'' も ``Wright Brothers'' も前例がある。

\paragraph{壁 2：funding がない.}
回避：DBR を自前で買う。これだけで既に「fund を持ち込む collaborator」になる。
¥500k は欧米の研究予算では微々たるもの。

\paragraph{壁 3：論文業績がない.}
回避：arXiv プレプリントが業績代わり。
具体的には Akkermans 解析の Python コードを GitHub 公開し、それと
関連付けた arXiv 論文 1 本があれば、最低限の信頼は得られる。

\paragraph{壁 4：visa / 訪問問題.}
回避：Zoom 会議で済ませる。実験は相手のラボでやってもらい、データを送って
もらう。実装フェーズだけ短期訪問（観光ビザの範囲内）。

\newpage

%======================================================================
\section{段階的実行プラン（独立研究者向け）}
%======================================================================

\subsection{Phase A：準備（1--2 ヶ月）}

\begin{enumerate}
\item Akkermans 公式の数値解を Python で実装、GitHub 公開
  \begin{itemize}
  \item リポジトリ名例：\texttt{akkermans-boyer-inverse-design}
  \item README で目的を明示、再現可能なノートブックを置く
  \end{itemize}
\item arXiv プレプリントを 1 本書く（6--8 ページ）
  \begin{itemize}
  \item タイトル：``Akkermans inverse design of artificial magnetic
    conductors for the experimental verification of Boyer's repulsive
    Casimir prediction''
  \item arXiv カテゴリ：quant-ph（または cond-mat.mes-hall）
  \end{itemize}
\item Layertec / Newport から DBR の見積もりを取る
  \begin{itemize}
  \item 設計仕様：本ガイド §3.3 と §7
  \item 見積書を PDF で保管
  \end{itemize}
\item 候補者上位 5 名の最近 3 年の論文を読破
\end{enumerate}

\subsection{Phase B：アプローチ（1--3 ヶ月）}

\begin{enumerate}
\item \textbf{Wave 1}：Iannuzzi にメール送信（候補 1）
  \begin{itemize}
  \item 1 ヶ月待つ
  \item 返信あり $\to$ オンライン会議
  \item 返信なし $\to$ Wave 2 へ
  \end{itemize}
\item \textbf{Wave 2}：Sushkov, Munday に同時メール（候補 2, 3）
  \begin{itemize}
  \item 並列で 1 ヶ月待つ
  \item 1 つでも返信あれば Wave 3 を保留
  \end{itemize}
\item \textbf{Wave 3}：Lamoreaux, Chan, 国内候補に並行メール
  \begin{itemize}
  \item ここまでで 3 ヶ月経過
  \item どこかから返信があれば共同研究合意へ
  \end{itemize}
\end{enumerate}

\subsection{Phase C：合意形成（1 ヶ月）}

\begin{enumerate}
\item オンライン会議で技術内容を擦り合わせ
\item 役割分担、論文オーサーシップ、データ所有権を明文化
\item 簡易 MoU（覚書）を作成・署名
\item DBR 発注、Layertec へ ¥400k 送金
\end{enumerate}

\subsection{Phase D：実験（2--3 ヶ月）}

\begin{enumerate}
\item DBR 到着、光学評価
\item 共同研究先で校正測定（Au-Au 標準）
\item DBR 測定本番
\item データ解析、誤差評価
\end{enumerate}

\subsection{Phase E：論文化（2--3 ヶ月）}

\begin{enumerate}
\item ドラフト（共同研究先 PI が筆頭著者）
\item 内部レビュー、再測定
\item Nature/PRL 投稿
\end{enumerate}

\textbf{合計期間：8--12 ヶ月}。

\newpage

%======================================================================
\section{代替プラン（全員に断られた場合）}
%======================================================================

\subsection{Plan B：DIY 化}

シナリオ B（自前構築、¥5M 級）に切り替え、独立研究者の自宅 / レンタル
スペースで実験を組む。

\textbf{ハードル}：
\begin{itemize}
\item AFM の自作は事実上不可能（精度が出ない）
\item MEMS チップの自作は MEMS ファブとの契約が必要
\item 振動絶縁・真空・温度制御の同時管理が必要
\end{itemize}

\textbf{現実的選択}：自作よりも \textbf{中古 AFM の購入}（eBay などで $20k--$50k 程度）。
JEOL や Park Systems の旧モデルなら入手可能。

\subsection{Plan C：別実験への切り替え}

Setup B' を諦め、より低コスト・より自前で可能な実験に切り替える：

\begin{itemize}
\item \textbf{Phase P}（Berry phase $\pi$ phononics、¥27k）：3D プリント結晶で
  $\pi$ Berry 位相のフォノニック実装。WB の既存ガイド
  \texttt{guide\_phase\_p\_comprehensive\_ja.tex} を参照。
\item \textbf{LC 回路実験}：単純な LC 共振器でエッジ状態を観測する
  ¥5k 級実験。\texttt{guide\_complete\_beginners\_ja.tex} 参照。
\end{itemize}

これらは Setup B' より直接性は低いが、\textbf{自前で完結} できる。

\subsection{Plan D：ブログ・arXiv ベースの公開議論}

実験は諦め、arXiv プレプリント + 個人ブログで提案を公開し、誰かが拾って
くれるのを待つ。これは長期的（年単位）戦略だが、完全に独立で進められる。

\newpage

%======================================================================
\section{テンプレート：初回メール}
%======================================================================

\subsection{Iannuzzi 向け（最有力候補用）}

\begin{tcolorbox}[colback=white,colframe=blue!50!black,title={英語メール例（カスタマイズ前）},breakable]
\small\ttfamily
Subject: A direct experimental extension of your 2004 PNAS paper:
Akkermans-designed DBR for Boyer's repulsive Casimir prediction\\

Dear Prof.\ Iannuzzi,\\

Your 2004 PNAS paper (Iannuzzi, Lisanti, Capasso) showed that the
Casimir force can be modulated by changing the optical state of one
mirror. This established the conceptual route I would like to extend.\\

Boyer's 1974 prediction that PEC--PMC boundaries give a repulsive
Casimir force has never been experimentally verified, because PMCs
do not exist as natural materials. Using the Akkermans--Dunne--Levy
formula (Optics of Aperiodic Structures, 2013, Ch.\ 10), we have
inverse-designed a SiO2/Ta2O5 DBR multilayer that implements the
DN-equivalent boundary condition over the relevant Casimir frequency
range. Predicted result: $F_{\rm DBR}/F_{\rm Au} = -7/8$.\\

Preprint with full derivation, transfer-matrix code, and design
specifications: arXiv:XXXX.YYYYY\\
GitHub (reproducible Python): github.com/wright-brothers/akkermans-boyer\\

I am an independent researcher and would like to propose a
collaboration where your group performs the measurement using your
existing AFM/MEMS Casimir setup, while I provide:\\
\hphantom{xx}- the DBR (already quoted at Layertec, EUR 3000, paid by me),\\
\hphantom{xx}- all theoretical/numerical analysis,\\
\hphantom{xx}- travel and labor for any in-lab work.\\

You and your group would be the lead authors on the resulting paper.
Even if the predicted sign reversal is not observed, the experiment
would be the first direct test of the Boyer 50-year prediction in
the PEC--PMC channel and is independently publishable in PRL or
Nature class.\\

Would you have 30 minutes for a Zoom call in the next two weeks?\\

Best regards,\\
Wright Brothers\\
Independent Researcher\\
[email address]\\
\end{tcolorbox}

\subsection{Sushkov 向け（CASPEr との接続版）}

ほぼ Iannuzzi 向けと同じテンプレート。違いは冒頭：

\begin{tcolorbox}[colback=white,colframe=blue!50!black,breakable]
\small\ttfamily
Subject: Boyer's 50-year unverified Casimir prediction --- Akkermans-designed
DBR test, complementing your CASPEr precision measurement expertise\\

Dear Prof.\ Sushkov,\\

Your work on CASPEr and the thermal Casimir force (Sushkov et al.\ 2011,
Nature Phys.) demonstrates exactly the kind of precision foundational
physics measurement that the proposed experiment requires.\\

Boyer's 1974 prediction... [以下、Iannuzzi 向けと同じ]
\end{tcolorbox}

\subsection{Munday 向け（Nature 2009 拡張版）}

\begin{tcolorbox}[colback=white,colframe=blue!50!black,breakable]
\small\ttfamily
Subject: Boyer's PEC--PMC Casimir prediction: a complementary route to
your 2009 Nature paper\\

Dear Prof.\ Munday,\\

Your 2009 Nature paper showed repulsive Casimir forces via dielectric
contrast in fluids. This is, to my knowledge, the only direct experimental
demonstration of Casimir repulsion. We propose to test the *other*
classical mechanism for Casimir repulsion --- Boyer's 1974 PEC--PMC
prediction, which has never been verified...\\
\end{tcolorbox}

\textbf{重要}：これらは「テンプレート」です。送信前に各候補の最近の論文を
読んで、相手の \emph{現在の研究} に接続する 1--2 文を追加すること。

\newpage

%======================================================================
\section{チェックリスト}
%======================================================================

\subsection{Phase A 準備チェックリスト}

\begin{itemize}
\item[$\square$] arXiv プレプリント完成、投稿済み
\item[$\square$] arXiv 番号を取得
\item[$\square$] GitHub リポジトリ公開
\item[$\square$] Akkermans Python コード再現性確認
\item[$\square$] Layertec から DBR 見積もり取得（PDF 保管）
\item[$\square$] 候補上位 5 名の連絡先確認
\item[$\square$] 候補上位 5 名の最近 3 年論文読破
\item[$\square$] 各候補向けにカスタマイズしたメールドラフト
\item[$\square$] DBR 購入資金 ¥500k の準備（または親類・支援者からの調達）
\item[$\square$] 訪問用パスポート有効期限確認
\end{itemize}

\subsection{Phase B アプローチチェックリスト}

\begin{itemize}
\item[$\square$] Wave 1（Iannuzzi）送信日の記録
\item[$\square$] 1 ヶ月後にフォローアップ送信、または Wave 2 へ
\item[$\square$] Wave 2（Sushkov, Munday）送信日の記録
\item[$\square$] Wave 3（Lamoreaux, Chan, 国内）送信日の記録
\item[$\square$] 全 send/receive を spreadsheet で追跡
\end{itemize}

\subsection{Phase C 合意形成チェックリスト}

\begin{itemize}
\item[$\square$] オンライン会議実施
\item[$\square$] 役割分担文書化
\item[$\square$] オーサーシップ合意
\item[$\square$] データ所有権明文化
\item[$\square$] MoU 署名
\item[$\square$] DBR 発注
\end{itemize}

%======================================================================
\section{まとめ}
%======================================================================

\begin{tcolorbox}[colback=green!5,colframe=green!50!black,title=本ガイドの一行まとめ]
Setup B' は \textbf{¥1M 程度・8--12 ヶ月}で実施可能な「Boyer 50 年予言の初検証」
実験である。独立研究者の最大の障害は装置でも資金でもなく、\textbf{適切な
共同研究先の確保}である。最有力候補は Davide Iannuzzi (Vrije Univ.\ Amsterdam
/ Optics11)。彼の 2004 PNAS 論文の自然な拡張として位置付け、自費で DBR を
持ち込み、論文の筆頭著者を彼に渡す形で提案すれば成功確率は 30--50\%。
``反重力'' は絶対に最初に出さない。
\end{tcolorbox}

\vspace{2em}

\textbf{関連文書}：
\begin{itemize}
\item \texttt{experimental\_guide\_antigravity\_ja.pdf}：6 段階の実験ロードマップ
\item \texttt{notes\_2026\_04\_10\_phase2\_audit\_ja.pdf}：Akkermans 解析の数値詳細
\item \texttt{topological\_antigravity.tex}：WB 反重力理論の数学的定式化
\item \texttt{boyer\_reinterpretation\_ja.tex}：Boyer 効果の数論的再解釈
\end{itemize}

\vspace{2em}

\textbf{重要な注意（再掲）}：本ガイドの候補者プロフィールは公開情報からの
主観的推測である。各人物の実際の応答性は接触するまで分からない。
本ガイドは順序付けの参考に過ぎず、最終判断は読者自身の調査と直接接触に
基づくべきである。

\end{document}
